\documentclass[12pt]{article}
\usepackage[utf8]{inputenc}
\usepackage[T1]{fontenc}
\usepackage{hyperref}
\title{Umbrian Language Resources for K-12 Curriculum Development}
\author{Compiled by Claude Sonnet 5.5 (claude-sonnet-5-5) for Andr\'as Kornai}
\date{}
\begin{document}
\maketitle
\section{Introduction}
Umbrian (ISO 639-3 xum) is an extinct Italic language of ancient Umbria, known from about 30 inscriptions dated from the 7th to the 1st century BC. Nearly all knowledge comes from the Iguvine Tables, seven bronze tablets found at Gubbio. There are no native speakers and no school materials were found in the search results.
\section{Primary Educational Resources}
No K-12 textbooks or primers were found in the search results.
\section{Linguistic and Reference Materials}
\begin{itemize}
\item Wikipedia, Umbrian language: \url{https://en.wikipedia.org/wiki/Umbrian_language}
\item Wikipedia, Iguvine Tablets, which mentions Poultney's 1959 \emph{The Bronze Tables of Iguvium} with English translation and glossary, and Devoto's 1948 edition: \url{https://en.wikipedia.org/wiki/Iguvine_Tablets}
\item Britannica, Umbrian language: \url{https://britannica.com/topic/Umbrian-language}
\item University of Florence ``Italia antica'' project page, The Umbrian language: \url{https://www.prin-italia-antica.unifi.it/p151.html}
\item Encyclopedia.pub entry: \url{https://encyclopedia.pub/entry/31388}
\end{itemize}
\section{Religious and Community Resources}
The Iguvine Tables record the rites of the Atiedian Brethren, twelve priests of Jupiter at Iguvium (modern Gubbio). No community organisations were identified in the search results.
\section{Digital and Audio Resources}
No audio or software resources were found.
\section{Curriculum Development Context}
Umbrian cannot be taught as a spoken language. Realistic school use is as a unit in Latin, classics or history: the Italic languages, the two scripts of the tablets (native Umbrian alphabet and Latin), and Gubbio's ancient rites. The vocabulary is poorly understood, so classroom materials must be modest. Printed editions such as Poultney's should be checked at libraries.
\end{document}